% Informe de cierre Sprint 2 - CuscoNodes
\documentclass[12pt,a4paper]{article}
\usepackage[utf8]{inputenc}
\usepackage[T1]{fontenc}
\usepackage[spanish]{babel}
\usepackage{graphicx}
\usepackage{hyperref}
\usepackage{array}
\usepackage{longtable}
\usepackage{caption}
\usepackage{geometry}
\usepackage{setspace}
\usepackage{booktabs}
\geometry{left=2.5cm,right=2.5cm,top=2.5cm,bottom=2.5cm}
\setlength{\parskip}{0.5em}
\setlength{\parindent}{0pt}
\renewcommand{\familydefault}{\rmdefault}

% Color y tamaño global: todo en negro
\usepackage{xcolor}
\color{black}

% Comandos para tamaños solicitados
\newcommand{\sectiontitle}[1]{\vspace{0.6cm}\noindent{\fontsize{13}{15}\selectfont\textbf{#1}}\par\vspace{0.15cm}\noindent{\fontsize{12}{14}\selectfont}}
\newcommand{\bodytext}[1]{{\fontsize{12}{14}\selectfont #1}\par}
\begin{document}

% Portada adaptada (en blanco y negro)
\begin{titlepage}
  \centering
  {\bfseries\Large UNIVERSIDAD ANDINA DEL CUSCO}\\[0.5cm]
  {\large FACULTAD DE INGENIERÍA Y ARQUITECTURA}\\[0.2cm]
  {\large ESCUELA PROFESIONAL DE INGENIERÍA DE SISTEMAS}\\[0.8cm]

  % Logo (usar logo en escala de grises en la misma carpeta si disponible)
  \includegraphics[width=0.35\textwidth]{logo_uac.png}\\[0.6cm]

  {\bfseries\Large INFORME DE CIERRE DE SPRINT}\\[0.4cm]
  {\bfseries\Large CuscoNodes - Sprint 2: Módulo de Percepción}\\[1.0cm]

  \begin{flushleft}
    \textbf{ASIGNATURA:}\hspace{1.5cm} Desarrollo de Proyecto / Ingeniería de Sistemas\\[0.3cm]
    \textbf{DOCENTE:}\hspace{2.3cm} Mg. Ing. Hugo Espetia Huamanga\\[0.3cm]
    \textbf{GRUPO:}\hspace{2.8cm} HABA\\[0.8cm]
    \textbf{INTEGRANTES:}\\[0.2cm]
    \hspace{3cm} Morales Ttito Honorio \hfill 100\\%
    \hspace{3cm} Torres Cáceres Sayo Michael \hfill 100\\%
    \hspace{3cm} Zavaleta Luna Sebastian (SM) \hfill 100\\%\\
  \end{flushleft}
  \vspace{1.5cm}

  \href{https://github.com/Honorio-Morales/cusconodes}{https://github.com/Honorio-Morales/cusconodes}\\[0.8cm]

  \vfill
  {\large CUSCO -- PERÚ}\\[0.3cm]
  {\large 2026}
\end{titlepage}

% Índice
{\fontsize{12}{14}\selectfont\tableofcontents}\newpage

% Introducción
\sectiontitle{Introducción}
\bodytext{Este documento formaliza la culminación del Sprint 2 del proyecto \textit{CuscoNodes}. Presenta los objetivos, alcance planificado, trabajo realizado, resultados obtenidos, evidencias y métricas que permiten justificar la aceptación del sprint y preparar la transición al siguiente ciclo de trabajo. Todo el formato del documento utiliza tipografía a tamaño 12 para el cuerpo y 13 para títulos, en escala de grises/negro y blanco.}

% Objetivos
\sectiontitle{Objetivos}
\bodytext{El objetivo principal de este sprint fue diseñar, implementar y validar el Módulo de Percepción (capa de extracción y pre-procesamiento de datos) para CuscoNodes, que incluye:}
\begin{itemize}
  \item Implementar scrapers para fuentes relevantes (RPP, PeruRail, El Comercio).
  \item Detectar y categorizar suspensiones y alertas críticas (PeruRailAnnouncements).
  \item Implementar un sistema de filtrado y scoring de criticidad orientado al turista (AlertFilter).
  \item Orquestar el pipeline completo (CuscoNodesScheduler) y preparar ejecución periódica (APScheduler).
  \item Documentar y almacenar salidas en formato JSON con metadatos.
\end{itemize}}

% Desarrollo del Sprint
\sectiontitle{Desarrollo del Sprint}
\bodytext{En este apartado se documenta con mayor detalle qué se hizo, por qué se eligieron las soluciones técnicas, cómo se validaron los resultados y cuáles fueron las decisiones operativas importantes. El equipo trabajó con un enfoque ágil, con reuniones de seguimiento cortas y revisiones técnicas al final del sprint.}

\subsection*{2. Objetivo del Sprint}
\bodytext{El Sprint 2 tuvo por objetivo entregar un módulo de Percepción robusto y verificable que proporcionara datos relevantes y filtrados para el sistema CuscoNodes. Específicamente:}
\begin{itemize}
  \item Desarrollar y validar scrapers para las fuentes seleccionadas (RPP, PeruRail, El Comercio).
  \item Detectar y clasificar interrupciones relevantes (suspensiones de trenes) con un scraper dedicado.
  \item Implementar un sistema de filtrado y scoring (AlertFilter) orientado a impactos turísticos.
  \item Orquestar el pipeline y asegurar persistencia y trazabilidad de los datos (nombres de archivo con timestamp, logs).
  \item Documentar y generar evidencia para cierre de sprint (commits, JSON de salida, capturas de Jira).
\end{itemize}

\subsection*{3. Alcance Planificado}
\bodytext{A continuación se listan las historias planificadas y su estado final. Cada historia está vinculada a entregables técnicos y artefactos de evidencia.}

\begin{longtable}{p{1.5cm} p{8cm} p{2.5cm} p{2.0cm}}
\toprule
\textbf{ID} & \textbf{Historia} & \textbf{Prioridad} & \textbf{Estado} \\
\midrule
CUN-12 & Configuración técnica: repositorio, entorno, estructura del proyecto & Alta & DONE \\
CUN-13 & RPPScraper: extracción de noticias de Cusco & Alta & DONE \\
CUN-14 & PeruRailScraper: extracción de horarios y PeruRailAnnouncementScraper: detección de suspensiones & Alta & DONE \\
CUN-15 & Almacenamiento local JSON con metadatos y archivos raw/processed & Media & DONE \\
CUN-16 & AlertFilter: filtrado, deduplicación y scoring de criticidad & Alta & DONE \\
CUN-17 & CuscoNodesScheduler: orquestador del pipeline y preparación para scheduling periódico & Alta & DONE \\
\bottomrule
\end{longtable}

\paragraph{Decisiones técnicas relevantes}
\bodytext{Se optó por una arquitectura modular: scrapers independientes por fuente, un módulo de filtrado separado y un orquestador que coordina ejecución, almacenamiento y generación de resúmenes. Esta separación facilita pruebas unitarias y despliegue gradual. Se eligieron JSONs locales para validación y pruebas; en producción se recomienda migrar a un storage versionado.}

% Resultados
\sectiontitle{Resultados}
\bodytext{Descripción de entregables y resultados obtenidos:}

\subsection*{4. Resultados Obtenidos}
\begin{longtable}{p{1.5cm} p{8cm} p{4cm}}
\toprule
\textbf{ID} & \textbf{Entregado} & \textbf{Evidencia} \\
\midrule
CUN-12 & Repositorio estructurado, configuración inicial & Commit: 22d4c30 en GitHub - https://github.com/Honorio-Morales/cusconodes \\
CUN-13 & RPPScraper: 15 artículos validados & data/raw/alertas_rpp_20260424_1922.json \\
CUN-14 & PeruRailScraper: 60+ horarios; PeruRailAnnouncementScraper & data/raw/horarios_perurail_20260424_1922.json; src/scrapers/perurail_announcement_scraper.py \\
CUN-15 & JSON storage + metadata & /data/processed/alertas_criticas_*.json \\
CUN-16 & AlertFilter implementado y validado & src/scrapers/alert_filter.py; resumen generado \\
CUN-17 & Scheduler y pipeline integradas & src/scheduler.py; ejecución local verificada \\
\bottomrule
\end{longtable}

% Evidencias
\sectiontitle{Evidencias}
\bodytext{A continuación se indican las evidencias generadas durante el sprint. Estas evidencias están almacenadas en el repositorio y en la carpeta `data/` para su revisión. Se recomienda incorporar capturas de Jira antes de la entrega final.}

\begin{figure}[h]
  \centering
  \includegraphics[width=0.8\textwidth]{jira_tablero.png}
  \caption*{Figura 1: Captura del tablero del Sprint en Jira (placeholder).}
\end{figure}

\bodytext{Evidencias técnicas y artefactos:}
\begin{itemize}
  \item Repositorio y commits: commit final `22d4c30` disponible en `main` y tag `sprint-2-percepcion`.
  \item Archivos raw: `data/raw/alertas_rpp_20260424_1922.json`, `data/raw/horarios_perurail_20260424_1922.json`, `data/raw/alertas_climaticas_20260424_1922.json`.
  \item Archivos procesados: `data/processed/alertas_criticas_*.json` y `resumen_alertas_*.json` (generados por `CuscoNodesScheduler`).
  \item Código: `src/scrapers/rpp_scraper.py`, `src/scrapers/perurail_scraper.py`, `src/scrapers/perurail_announcement_scraper.py`, `src/scrapers/alert_filter.py`, `src/scheduler.py`.
  \item Documentación: `docs/TAREAS_JIRA.md`, `docs/SPRINT2_FINAL.md`, `docs/INFORME_SPRINT2.tex`.
  \item Capturas de Jira (por adjuntar): `jira_CUN-12.png`, ..., `jira_CUN-17.png`.
\end{itemize}

\paragraph{Validación realizada}
\bodytext{Para validar la efectividad del módulo de percepción se realizaron las siguientes acciones:}
\begin{itemize}
  \item Ejecución del pipeline completo en ambiente local para verificar integridad de parseo y consistencia de campos.
  \item Validación manual de una muestra de 30 artículos/historias para comprobar normalización de campos (`titulo`, `enlace`, `fuente`, `tipo`, `fecha_scrape`).
  \item Comparación de conteos: 15 artículos RPP, 60+ registros PeruRail, 200+ registros climáticos (total 345+ registros integrados).
  \item Revisión de casos detectados por `PeruRailAnnouncementScraper` para verificar que las suspensiones fueran correctamente clasificadas y con urgencia calculada.
\end{itemize}

% Incidencias y Riesgos
\sectiontitle{Incidencias y Riesgos}
\begin{longtable}{p{6cm} p{3.5cm} p{5cm}}
\toprule
\textbf{Problema} & \textbf{Impacto} & \textbf{Acción} \\
\midrule
Inconsistencia en campos de JSON entre scrapers & Alto & Normalizar campos: `titulo`, `enlace`, `fuente`, `tipo`, `fecha_scrape`. Actualización de parse_article() en scrapers. \\
Dependencia del sitio PeruRail para detección de suspensiones & Medio-Alto & Implementar heurística de palabras claves y manejo de errores; separar scrapers por tipo de contenido. \\
Archivos de datos grandes (alertas climáticas) ignorados por .gitignore & Bajo-Medio & Forzar inclusión con `git add -f` para datos de validación; en producción usar almacenamiento externo o dataset versionado. \\
\bottomrule
\end{longtable}

% Métricas
\sectiontitle{Métricas}
\bodytext{A continuación métricas relevantes del Sprint 2:}

\begin{itemize}
  \item Velocidad (Velocity): 16 story points completados
  \item Tareas completadas: 6/6
  \item Bugs detectados: 0 críticos en producción (problemas menores documentados)
  \item Tiempo real vs estimado: En promedio, las tareas se completaron dentro de +/- 10% del tiempo estimado. (Detalle en Anexos)
\end{itemize}

% Conclusiones
\sectiontitle{Conclusiones}
\bodytext{El Sprint 2 alcanzó sus objetivos principales. Resumen de conclusiones clave:}
\begin{itemize}
  \item El Módulo de Percepción entrega datos estructurados y trazables desde múltiples fuentes, lo que habilita las capas superiores (motor de razonamiento y notificaciones).
  \item El sistema distingue entre contenido informativo y alertas de impacto turístico mediante reglas y scoring, reduciendo el ruido para usuarios finales.
  \item La detección de suspensiones en PeruRail está operativa y permite priorizar eventos por urgencia.
  \item La arquitectura modular facilita extensión (nuevos scrapers, integración de IA) y despliegue.
\end{itemize}
\bodytext{Limitaciones detectadas:}
\begin{itemize}
  \item Dependencia de cambios en la estructura HTML de fuentes externas; se mitigó con manejo de excepciones y tests de regresión simples.
  \item Almacenamiento local adecuado para pruebas, pero se recomienda migración a storage versionado para producción y control de datos sensibles.
\end{itemize}

% Recomendaciones
\sectiontitle{Recomendaciones}
\bodytext{Para el siguiente sprint (Sprint 3) se recomienda:}
\begin{itemize}
  \item Implementar el Motor de Razonamiento (IA) para clasificación automática y resumen de alertas.
  \item Añadir traducción automática para multilanguage (EN, FR, PT).
  \item Implementar pipeline de notificaciones (WhatsApp, Email, Push).
  \item Mover almacenamiento de datos de validación a un storage versionado o S3 para producción.
  \item Automatizar compilación del informe y generación de PDF con las capturas de Jira incluidas.
\end{itemize}

% Anexos
\sectiontitle{Anexos}
\bodytext{Archivos y recursos anexos (adjuntar antes de entregar):}
\begin{itemize}
  \item `docs/TAREAS_JIRA.md` (descripciones completas de las tareas)
  \item `docs/SPRINT2_FINAL.md` (cierre formal completo)
  \item `data/raw/alertas_rpp_20260424_1922.json`
  \item `data/raw/horarios_perurail_20260424_1922.json`
  \item `data/raw/alertas_climaticas_20260424_1922.json`
  \item Capturas Jira: `jira_tablero.png`, `jira_CUN-12.png`, ..., `jira_CUN-17.png` (adjuntar antes de compilar)
  \item Registro de commits relevantes: `git log --oneline` (ver `22d4c30`)
\end{itemize}

\vspace{1cm}
\bodytext{Documento generado: 27 de Abril de 2026.}

\end{document}
